% GENERATED FILE. Edit canonical lessons, then run npm run generate:books.
% canonical-source-hash: fnv1a64:275e6df4a5ccfa21
% canonical-lessons: GE-C330-standpunkt, GE-C330-ansicht, GE-C330-auffassung, GE-C330-gesichtspunkt, GE-C330-vertreten

\chapter{Mein Standpunkt}
\label{ch:ge-mein-standpunkt}
\hlchaptermodality{330}

\begin{chapteropening}
\textbf{By the end of this chapter:} \emph{I can state my point of view, say how I see a question and say which view I stand up for.}

Complete the last lesson of chapter 330: I can state my point of view, say how I see a question and say which view I stand up for.
\end{chapteropening}

\section[der Standpunkt]{der Standpunkt — the standpoint, the point of view}

\label{lesson:GE-C330-standpunkt}



\begin{quote}
Before the new one: say the German for at least, then the German for during.
\end{quote}

\subsection*{You'll want to know: der Standpunkt}
\textbf{der Standpunkt} — \textquotedblleft{}the standpoint, the point of view\textquotedblright{}.

\begin{quote}
Von meinem Standpunkt aus ist das zu teuer.
\end{quote}

\begin{quote}
\emph{From my point of view, that is too expensive.}
\end{quote}

Plural: \emph{die Standpunkte}. \emph{Von meinem Standpunkt aus} — from where I stand.

\begin{cousinweb}
\emph{Standpunkt} is built from \emph{Stand} \textquotedblleft{}standing, position\textquotedblright{}, inherited from Old High German \emph{stand} and Proto-Germanic \emph{standaz} (EtymDB), of the family of \emph{stehen} \textquotedblleft{}to stand\textquotedblright{}, and \emph{Punkt}, borrowed from Latin \emph{punctum} \textquotedblleft{}a prick, a point\textquotedblright{} (EtymDB; Lewis \& Short): the spot you stand on.
\end{cousinweb}

\begin{grammarlens}[title={What you've built}]
One more word for telling, explaining and describing at B1.
\end{grammarlens}

\subsection*{Your turn}
\begin{itemize}
\raggedright
  \item \emph{Say it:} \emph{der Standpunkt}
  \item \emph{Say it:} \emph{der Standpunkt}, once more
  \item \emph{Say it:} a sentence of your own with \emph{der Standpunkt}
  \item \emph{Recall:} say the German for at least, then the German for during, then say \emph{der Standpunkt} again
\end{itemize}

\begin{tcolorbox}[breakable,colback=teal!4,colframe=teal!35!black,title={Before you move on}]
What does der Standpunkt mean? (\textquotedblleft{}The standpoint, the point of view\textquotedblright{}.) Say it once more.
\end{tcolorbox}
\section[die Ansicht]{die Ansicht — the view, the opinion}

\label{lesson:GE-C330-ansicht}



\begin{quote}
Before the new one: say the German for during, then the German for the standpoint.
\end{quote}

\subsection*{You'll want to know: die Ansicht}
\textbf{die Ansicht} — \textquotedblleft{}the view, the opinion\textquotedblright{}.

\begin{quote}
Meiner Ansicht nach ist der Film zu lang.
\end{quote}

\begin{quote}
\emph{In my view, the film is too long.}
\end{quote}

\emph{Meiner Ansicht nach} works like \emph{meiner Meinung nach}. Plural: \emph{die Ansichten}; \emph{die Ansichtskarte} is a picture postcard.

\begin{cousinweb}
\emph{Ansicht} is formed on \emph{ansehen} \textquotedblleft{}to view, to regard\textquotedblright{}, from \emph{an} and \emph{sehen} \textquotedblleft{}to see\textquotedblright{}, which is inherited from Old High German \emph{sehan} and Proto-Germanic \emph{sehwaną} (EtymDB): the way a thing looks to you.
\end{cousinweb}

\begin{grammarlens}[title={What you've built}]
One more word for telling, explaining and describing at B1.
\end{grammarlens}

\subsection*{Your turn}
\begin{itemize}
\raggedright
  \item \emph{Say it:} \emph{die Ansicht}
  \item \emph{Say it:} \emph{die Ansicht}, once more
  \item \emph{Say it:} a sentence of your own with \emph{die Ansicht}
  \item \emph{Recall:} say the German for during, then the German for the standpoint, then say \emph{die Ansicht} again
\end{itemize}

\begin{tcolorbox}[breakable,colback=teal!4,colframe=teal!35!black,title={Before you move on}]
What does die Ansicht mean? (\textquotedblleft{}The view, the opinion\textquotedblright{}.) Say it once more.
\end{tcolorbox}
\section[die Auffassung]{die Auffassung — the understanding, the conception (of a question)}

\label{lesson:GE-C330-auffassung}



\begin{quote}
Before the new one: say the German for the standpoint, then the German for the view.
\end{quote}

\subsection*{You'll want to know: die Auffassung}
\textbf{die Auffassung} — \textquotedblleft{}the understanding, the conception (of a question)\textquotedblright{}.

\begin{quote}
Nach meiner Auffassung muss die Stadt mehr Busse kaufen.
\end{quote}

\begin{quote}
\emph{As I understand it, the town has to buy more buses.}
\end{quote}

More formal than \emph{die Meinung}. \emph{Ich bin der Auffassung, dass …} — I take the view that …

\begin{cousinweb}
\emph{Auffassung} is the noun of \emph{auffassen} \textquotedblleft{}to take in, to grasp\textquotedblright{}, built on \emph{fassen} \textquotedblleft{}to seize, to hold\textquotedblright{}, inherited from Old High German \emph{fazzōn} and Proto-Germanic \emph{fatōną} (EtymDB): the way you take hold of an idea.
\end{cousinweb}

\begin{grammarlens}[title={What you've built}]
One more word for telling, explaining and describing at B1.
\end{grammarlens}

\subsection*{Your turn}
\begin{itemize}
\raggedright
  \item \emph{Say it:} \emph{die Auffassung}
  \item \emph{Say it:} \emph{die Auffassung}, once more
  \item \emph{Say it:} a sentence of your own with \emph{die Auffassung}
  \item \emph{Recall:} say the German for the standpoint, then the German for the view, then say \emph{die Auffassung} again
\end{itemize}

\begin{tcolorbox}[breakable,colback=teal!4,colframe=teal!35!black,title={Before you move on}]
What does die Auffassung mean? (\textquotedblleft{}The understanding, the conception (of a question)\textquotedblright{}.) Say it once more.
\end{tcolorbox}
\section[der Gesichtspunkt]{der Gesichtspunkt — the aspect, the angle (of a question)}

\label{lesson:GE-C330-gesichtspunkt}



\begin{quote}
Before the new one: say the German for the view, then the German for the understanding.
\end{quote}

\subsection*{You'll want to know: der Gesichtspunkt}
\textbf{der Gesichtspunkt} — \textquotedblleft{}the aspect, the angle (of a question)\textquotedblright{}.

\begin{quote}
Unter diesem Gesichtspunkt ist die Idee gar nicht so schlecht.
\end{quote}

\begin{quote}
\emph{From this angle, the idea is not so bad at all.}
\end{quote}

Plural: \emph{die Gesichtspunkte}. \emph{Gesicht} here keeps its old sense, sight.

\begin{cousinweb}
\emph{Gesichtspunkt} joins \emph{Gesicht}, here in its older sense \textquotedblleft{}sight\textquotedblright{}, from the family of \emph{sehen} \textquotedblleft{}to see\textquotedblright{}, to \emph{Punkt} \textquotedblleft{}point\textquotedblright{}, from Latin \emph{punctum} (EtymDB). It is usually explained as a loan translation of the Latin term of optics \emph{punctum visus}, \textquotedblleft{}the point of sight\textquotedblright{}.
\end{cousinweb}

\begin{grammarlens}[title={What you've built}]
One more word for telling, explaining and describing at B1.
\end{grammarlens}

\subsection*{Your turn}
\begin{itemize}
\raggedright
  \item \emph{Say it:} \emph{der Gesichtspunkt}
  \item \emph{Say it:} \emph{der Gesichtspunkt}, once more
  \item \emph{Say it:} a sentence of your own with \emph{der Gesichtspunkt}
  \item \emph{Recall:} say the German for the view, then the German for the understanding, then say \emph{der Gesichtspunkt} again
\end{itemize}

\begin{tcolorbox}[breakable,colback=teal!4,colframe=teal!35!black,title={Before you move on}]
What does der Gesichtspunkt mean? (\textquotedblleft{}The aspect, the angle (of a question)\textquotedblright{}.) Say it once more.
\end{tcolorbox}
\section[vertreten]{vertreten — to advocate (a view), to represent; to stand in for}

\label{lesson:GE-C330-vertreten}



\begin{quote}
Before the new one: say the German for the understanding, then the German for the aspect.
\end{quote}

\subsection*{You'll want to know: vertreten}
\textbf{vertreten} — \textquotedblleft{}to advocate (a view), to represent; to stand in for\textquotedblright{}.

\begin{quote}
Sie vertritt die Meinung, dass Kinder weniger fernsehen sollten.
\end{quote}

\begin{quote}
\emph{She holds the view that children should watch less television.}
\end{quote}

Inseparable and strong: \emph{sie vertritt, sie hat vertreten}. At work, \emph{jemanden vertreten} is to stand in for someone.

\begin{cousinweb}
\emph{vertreten} is built from \emph{ver-} and \emph{treten} \textquotedblleft{}to step, to tread\textquotedblright{}, inherited from Old High German \emph{tretan} and Proto-Germanic \emph{tredaną} (EtymDB): you step into someone's place, or step forward for a view.
\end{cousinweb}

\begin{grammarlens}[title={What you've built}]
One more word for telling, explaining and describing at B1.
\end{grammarlens}

\subsection*{Your turn}
\begin{itemize}
\raggedright
  \item \emph{Say it:} \emph{vertreten}
  \item \emph{Say it:} \emph{vertreten}, once more
  \item \emph{Say it:} a sentence of your own with \emph{vertreten}
  \item \emph{Recall:} say the German for the understanding, then the German for the aspect, then say \emph{vertreten} again
\end{itemize}

\begin{tcolorbox}[breakable,colback=teal!4,colframe=teal!35!black,title={Before you move on}]
What does vertreten mean? (\textquotedblleft{}To advocate (a view), to represent; to stand in for\textquotedblright{}.) Say it once more.
\end{tcolorbox}
